\begin{helper}[Transporting a good witness to an extended endpoint]
Let $K_J$ be a good hierarchy on a corresponding ETT over $J$, and let
$K_U$ be a hierarchy on a corresponding ETT over $U$, where both ETTs
correspond to the same original ETT and coloring $\sigma$.  Assume the
bases have the same core and closure, the divider counts agree, and under
the canonical index identification all nonterminal stages, starred stages,
indexed stored sets, and stage-missing sets agree.  Suppose also that
$J$ reaches $U$ by a $\sigma$-Tashkinov augmentation sequence.  If every
last-index Gamma pair of $K_J$ avoids the restricted segment colors at
$U$, then $K_U$ has a good witness whose Gamma pairs at every index equal
the corresponding Gamma pairs of $K_J$.
\end{helper}
\begin{proof}
Define the Gamma pairs for $K_U$ by copying those of $K_J$ through the
canonical divider-index identification.  The two-color cardinality,
pairwise-disjointness, and unique-refresh fields of
\noderef{GoodHierarchy} transport immediately across the assumed
equalities of indexed stored sets, stage-missing sets, starred stages, and
Gamma pairs.  The internal minus-closed fields use only the copied stages,
stored sets, and Gamma pairs and therefore transport in the same way.  For
the PE-base field, the two
\noderef{CorrespondingExtendedTashkinovTree} certificates identify the
terminal extension kind and terminal ladder tree with the same original
data, while the assumed core and closure equalities identify the remaining
base data; hence the old PE-base implication rewrites to the new one.

It remains to prove avoidance, because the segment-color definition also
uses the entire terminal tree.  By
\noderef{TashkinovAugmentationSequenceInitialSegment}, the augmentation
sequence gives an exact ordered prefix $J\preccurlyeq U$.  Fix an index
$j<d$, a stored color $\eta$, and let $B$ be the common copied next stage.
Iterating the nesting fields of \noderef{TashkinovHierarchy} gives
$B\preccurlyeq J$, hence $|B|\leq |J|$.  Apply
\noderef{HierarchyFirstMissingPositionPrefixCutoff} to
$J\preccurlyeq U$ with cutoff $|B|$.  It yields
\[
 \min\{p_J(\eta),|B|\}=\min\{p_U(\eta),|B|\}.
\]
This equality includes the case in which no vertex misses $\eta$, because
the cited theorem handles the totalized last-position fallback.

Unfold \noderef{HierarchySegmentColors}.  The lower starred stage and the
edge map of $B$ agree by hypothesis.  At every edge index of $B$, the
lower-bound test agrees, the displayed cutoff equality identifies the
strict upper-bound test, and the edge-color test is literally the same.
Thus the old and new segment-color sets are equal.  The old avoidance
field of the good witness on $K_J$ consequently supplies avoidance for
every $j<d$.

At $j=d$, the new next stage is $U$, the starred stage is the common last
starred tree, and \noderef{HierarchySegmentColors} is exactly the
restricted filter in the hypothesis.  Equality of indexed stored and
stage-missing sets transports the containment part of the old avoidance
field, while the endpoint hypothesis supplies disjointness from the new
segment set.  Hence avoidance holds at the last index as well.  All fields
of \noderef{GoodHierarchy} are now verified, and the construction gives
the asserted Gamma equalities.
\end{proof}
